\begin{abstract}
Generative world models for autonomous driving (AD) have become a trending topic. 
Unlike the widely studied image modality, in this work we explore generative world models for LiDAR data.
Existing generation methods for LiDAR data only support single frame generation, while existing prediction approaches require multiple frames of historical input and can only deterministically predict multiple frames at once, lacking interactivity.
Both paradigms fail to support long-horizon interactive generation.
%To this end, we introduce \textbf{LaGen}, a framework capable of autoregressively generating long-horizon LiDAR scenes frame by frame. 
%LaGen can generate high-fidelity scene point clouds conditioned on a single-frame start LiDAR input and utilizes bounding box conditions.
To this end, we introduce \textbf{LaGen}, which to the best of our knowledge is the first framework capable of frame-by-frame autoregressive generation of long-horizon LiDAR scenes.
LaGen is able to take a single-frame LiDAR input as a starting point and effectively utilize bounding box information as conditions to generate high-fidelity 4D scene point clouds.
In addition, we introduce a scene decoupling estimation module to enhance the model's interactive generation capability for object-level content, as well as a noise modulation module to mitigate error accumulation during long-horizon generation.
We construct a protocol based on nuScenes for evaluating long-horizon LiDAR scene generation. 
Experimental results comprehensively demonstrate LaGen outperforms state-of-the-art LiDAR generation and prediction models, especially on the later frames.

\vspace{-4mm}

\end{abstract}